\documentclass{beamer}

\usepackage[T2A]{fontenc}
\usepackage[utf8]{inputenc}
\usepackage[english,russian]{babel}
\pdfmapfile{+sansmathaccent.map}


\mode<presentation>
{
	\usetheme{Warsaw} % or try Darmstadt, Madrid, Warsaw, Rochester, CambridgeUS, ...
	\usecolortheme{seahorse} % or try seahorse, beaver, crane, wolverine, ...
	\usefonttheme{serif}  % or try serif, structurebold, ...
	\setbeamertemplate{navigation symbols}{}
	\setbeamertemplate{caption}[numbered]
} 


%%%%%%%%%%%%%%%%%%%%%%%%%%%%
% itemize settings


%%%%%%%%%%%%%%%%%%%%%%%%%%%%
% itemize settings

\definecolor{myhotpink}{RGB}{255, 80, 200}
\definecolor{mywarmpink}{RGB}{255, 60, 160}
\definecolor{mylightpink}{RGB}{255, 80, 200}
\definecolor{mypink}{RGB}{255, 30, 80}
\definecolor{mydarkpink}{RGB}{155, 25, 60}

\definecolor{mypaleblue}{RGB}{240, 240, 255}
\definecolor{mylightblue}{RGB}{120, 150, 255}
\definecolor{myblue}{RGB}{90, 90, 255}
\definecolor{mygblue}{RGB}{70, 110, 240}
\definecolor{mycourseblue}{RGB}{120, 140, 255}
\definecolor{mycourselightblue}{RGB}{245, 250, 255}
\definecolor{mydarkblue}{RGB}{0, 0, 180}
\definecolor{myblackblue}{RGB}{40, 40, 120}

\definecolor{mygreen}{RGB}{0, 200, 0}
\definecolor{mygreen2}{RGB}{245, 255, 230}
\definecolor{mydarkgreen}{RGB}{0, 120, 0}
\definecolor{mydarkgreen}{RGB}{0, 120, 0}


\definecolor{mygray}{gray}{0.8}
\definecolor{mydarkgray}{RGB}{80, 80, 160}
\definecolor{mylightgray}{RGB}{160, 160, 160}

\definecolor{mydarkred}{RGB}{160, 30, 30}
\definecolor{mylightred}{RGB}{255, 150, 150}
\definecolor{myred}{RGB}{200, 110, 110}
\definecolor{myblackred}{RGB}{120, 40, 40}

\definecolor{mypink}{RGB}{255, 30, 80}
\definecolor{myhotpink}{RGB}{255, 80, 200}
\definecolor{mywarmpink}{RGB}{255, 60, 160}
\definecolor{mylightpink}{RGB}{255, 80, 200}
\definecolor{mydarkpink}{RGB}{155, 25, 60}
\definecolor{mywhitepink}{RGB}{255, 240, 240}

\definecolor{mydarkcolor}{RGB}{60, 25, 155}
\definecolor{mylightcolor}{RGB}{130, 180, 250}



\setbeamertemplate{itemize items}[default]

\setbeamertemplate{itemize item}{\color{myblackblue}$\blacksquare$}
\setbeamertemplate{itemize subitem}{\color{mygblue}$\blacktriangleright$}
\setbeamertemplate{itemize subsubitem}{\color{mygray}$\blacksquare$}

\setbeamercolor{palette quaternary}{fg=white,bg=mycourseblue}
\setbeamercolor{titlelike}{parent=palette quaternary}

\setbeamercolor{palette quaternary2}{fg=black,bg=mycourselightblue}
\setbeamercolor{frametitle}{parent=palette quaternary2}

\setbeamerfont{frametitle}{size=\Large,series=\scshape}
\setbeamerfont{framesubtitle}{size=\normalsize,series=\upshape}





%%%%%%%%%%%%%%%%%%%%%%%%%%%%
% block settings

\setbeamercolor{block title}{bg=red!30,fg=black}

\setbeamercolor*{block title example}{bg=mygreen!40!white,fg=black}

\setbeamercolor*{block body example}{fg= black, bg= mygreen2}


%%%%%%%%%%%%%%%%%%%%%%%%%%%%
% URL settings
\hypersetup{
	colorlinks=true,
	linkcolor=blue,
	filecolor=blue,      
	urlcolor=blue,
}

%%%%%%%%%%%%%%%%%%%%%%%%%%

\renewcommand{\familydefault}{\rmdefault}


\newcommand{\mycolor}[1] {
\setbeamercolor{normal text}{fg={#1}} % Recolors standard body text
\setbeamercolor{structure}{fg={#1}}   % Recolors bullets and structure items
\usebeamercolor[fg]{normal text}      % Applies the normal text color immediately
}


\usepackage{amsmath}
\usepackage{mathtools}

\usepackage{cancel}
\usepackage{centernot}

\usepackage{changepage}
\usepackage{subcaption}

\usepackage{qrcode}

\newcommand{\bo}[1] {\mathbf{#1}}
\newcommand{\R}{\mathbb{R}} 
\newcommand{\T}{^\top}     



\newcommand{\mydate}{Fall 2026}

\newcommand{\mygit}{\textcolor{blue}{\href{https://github.com/SergeiSa/Advanced-Control-2026}{github.com/SergeiSa/Advanced-Control-2026}}}

\newcommand{\myqr}{ \textcolor{black}{\qrcode[height=1.5in]{https://github.com/SergeiSa/Advanced-Control-2026}}
}




\newcommand{\bref}[2] {\textcolor{blue}{\href{#1}{#2}}}

%%%%%%%%%%%%%%%%%%%%%%%%%%%%
% code settings

\usepackage{listings}
\usepackage{color}
% \definecolor{mygreen}{rgb}{0,0.6,0}
% \definecolor{mygray}{rgb}{0.5,0.5,0.5}
\definecolor{mymauve}{rgb}{0.58,0,0.82}
\lstset{ 
	backgroundcolor=\color{white},   % choose the background color; you must add \usepackage{color} or \usepackage{xcolor}; should come as last argument
	basicstyle=\footnotesize,        % the size of the fonts that are used for the code
	breakatwhitespace=false,         % sets if automatic breaks should only happen at whitespace
	breaklines=true,                 % sets automatic line breaking
	captionpos=b,                    % sets the caption-position to bottom
	commentstyle=\color{mygreen},    % comment style
	deletekeywords={...},            % if you want to delete keywords from the given language
	escapeinside={\%*}{*)},          % if you want to add LaTeX within your code
	extendedchars=true,              % lets you use non-ASCII characters; for 8-bits encodings only, does not work with UTF-8
	firstnumber=0000,                % start line enumeration with line 0000
	frame=single,	                   % adds a frame around the code
	keepspaces=true,                 % keeps spaces in text, useful for keeping indentation of code (possibly needs columns=flexible)
	keywordstyle=\color{blue},       % keyword style
	language=Octave,                 % the language of the code
	morekeywords={*,...},            % if you want to add more keywords to the set
	numbers=left,                    % where to put the line-numbers; possible values are (none, left, right)
	numbersep=5pt,                   % how far the line-numbers are from the code
	numberstyle=\tiny\color{mygray}, % the style that is used for the line-numbers
	rulecolor=\color{black},         % if not set, the frame-color may be changed on line-breaks within not-black text (e.g. comments (green here))
	showspaces=false,                % show spaces everywhere adding particular underscores; it overrides 'showstringspaces'
	showstringspaces=false,          % underline spaces within strings only
	showtabs=false,                  % show tabs within strings adding particular underscores
	stepnumber=2,                    % the step between two line-numbers. If it's 1, each line will be numbered
	stringstyle=\color{mymauve},     % string literal style
	tabsize=2,	                   % sets default tabsize to 2 spaces
	title=\lstname                   % show the filename of files included with \lstinputlisting; also try caption instead of title
}


%%%%%%%%%%%%%%%%%%%%%%%%%%%%
% URL settings
\hypersetup{
	colorlinks=false,
	linkcolor=blue,
	filecolor=blue,      
	urlcolor=blue,
}

%%%%%%%%%%%%%%%%%%%%%%%%%%

%%%%%%%%%%%%%%%%%%%%%%%%%%%%
% tikz settings

\usepackage{tikz}
\tikzset{every picture/.style={line width=0.75pt}}
\usetikzlibrary{shapes.geometric, decorations.markings, positioning, patterns, calc, arrows.meta}

\newcommand{\myqrframe}{
	\begin{frame}
		\begin{center}
			\centerline{Слайды доступны на Github:}
			\bigskip
			\centerline{\mygit}
			\bigskip
			\myqr
		\end{center}
	\end{frame}
}


\subtitle{Современные методы управления роботами}
\author{Сергей Савин}
\date{2026}



\title{Q-обучение}
\begin{document}
	\maketitle
	
	
	\begin{frame}{Содержание}
		\begin{itemize}
			\item Обучение с подкреплением на основе модели и без модели
			\item Терминология: от оптимального управления к RL
			\item Возврат, награды и коэффициент дисконтирования
			\item Функция ценности действия $Q^\pi(s,a)$ и $Q^*(s,a)$
			\item Соотношение Беллмана и уравнение оптимальности
			\item Оптимальная политика из $Q^*$
			\item Временное разностное обновление (SARSA)
			\item Q-обучение (off-policy обновление)
			\item Глубокое Q-обучение (DQN)
			\item Литература
		\end{itemize}
	\end{frame}
	
	
	\begin{frame}{Обучение с подкреплением}
		\begin{flushleft}
			
			\textbf{Оптимальное управление на основе модели}
			
			Мы знаем модель системы:
			\[
			x_{k+1}=f(x_k,u_k,w_k).
			\]
			Это позволяет нам вычислять (или приближать) функции ценности
			и выводить оптимальную политику управления.
			
			\vspace{0.4cm}
			
			\textbf{Обучение с подкреплением без модели}
			
			Мы не знаем модель системы. Вместо этого мы собираем
			наблюдения и учимся непосредственно на данных.
			
			\vspace{0.3cm}
			
			\begin{itemize}
				\item Мы можем обучать функции ценности по наблюдаемым переходам.
				\item Как только мы знаем оптимальную функцию ценности, нахождение
				оптимальной политики может быть очень простым.
			\end{itemize}
		\end{flushleft}
	\end{frame}
	
	
	\begin{frame}{Терминология: связь с оптимальным управлением}
		
		\centering
		\begin{tabular}{ll}
			\textbf{Оптимальное управление} & \textbf{Обучение с подкреплением} \\
			\hline
			Система & Среда \\
			Контроллер & Агент \\
			Состояние $x$ & Состояние $s$ \\
			Управляющее воздействие $u$ & Действие $a$ \\
			Стоимость этапа $l(x,u)$ & Награда $r(s,a)$ \\
			Стоимость $J$ (минимизировать) & Возврат $J$ (максимизировать) \\
			Закон управления & Политика $\pi$
		\end{tabular}
		
		\vspace{0.4cm}
		
		\begin{itemize}
			\item Агент взаимодействует со средой.
			\item Агент наблюдает состояние и выбирает действие.
			\item Среда эволюционирует и выдаёт награду.
		\end{itemize}
		
	\end{frame}
	
	
	\begin{frame}{Функция возврата и награды}
		\begin{flushleft}
			Функция награды
			\[
			r_k = r(s_k,a_k)
			\]
			описывает немедленную награду, связанную с текущим
			состоянием и действием.
			
			\vspace{0.3cm}
			
			Дисконтированный возврат с шага $k$ равен
			\[
			G_k =
			\sum_{i=k}^{\infty}
			\gamma^{i-k}r(s_i,a_i).
			\]
			
			где $0<\gamma<1$ — коэффициент дисконтирования.
			
			\vspace{0.3cm}
			
			\begin{itemize}
				\item Немедленные награды имеют наибольший вес.
				\item Будущие награды прогрессивно дисконтируются.
				\item Цель — максимизировать ожидаемый возврат.
			\end{itemize}
			
		\end{flushleft}
	\end{frame}
	
	
	\begin{frame}{Функция ценности действия}
		
		\begin{flushleft}
			\textbf{On-policy функция ценности действия} — это
			\[
			Q^\pi(s,a)=
			\mathbb{E}_\pi\left[
			\sum_{k=0}^{\infty}\gamma^k r_{t+k}
			\;\middle|\;
			s_t=s,\ a_t=a
			\right].
			\]
			
			\vspace{0.2cm}
			
			Она описывает ожидаемый дисконтированный возврат, когда:
			\begin{itemize}
				\item Мы берём действие $a$ в состоянии $s$ сейчас.
				\item Мы следуем политике $\pi$ в дальнейшем.
			\end{itemize}
			
			\vspace{0.3cm}
			
			\textbf{Оптимальная функция ценности действия} — это
			\[
			Q^*(s,a)=\max_\pi Q^\pi(s,a).
			\]
			
			Она описывает наилучший ожидаемый дисконтированный возврат, достижимый
			при взятии действия $a$ в состоянии $s$ и следовании
			оптимальной политике в дальнейшем.
			
		\end{flushleft}
	\end{frame}
	
	
	\begin{frame}{Соотношение Беллмана}
		\begin{flushleft}
			
			Мы можем разложить дисконтированный возврат:
			\[
			\begin{aligned}
				G_k
				&=r_k+\gamma r_{k+1}+\gamma^2r_{k+2}+\ldots\\
				&=r_k+\gamma G_{k+1}.
			\end{aligned}
			\]
			
			\vspace{0.3cm}
			
			Взятие математического ожидания даёт уравнение Беллмана.
			
			\textbf{Оценка политики:}
			\[
			Q^\pi(s,a)=
			\mathbb{E}_\pi\left[
			r+\gamma Q^\pi(s',a')
			\right].
			\]
			
			\vspace{0.2cm}
			
			\textbf{Уравнение оптимальности:}
			\[
			Q^*(s,a)=
			\mathbb{E}\left[
			r+\gamma\max_{a'}Q^*(s',a')
			\right].
			\]
			
		\end{flushleft}
	\end{frame}
	
	
	\begin{frame}{Оптимальная политика}
		\begin{flushleft}
			
			Если мы знаем оптимальную функцию ценности действия,
			мы можем получить оптимальную политику, выбирая
			действие с наибольшим $Q$-значением.
			
			\[
			\boxed{
				\pi^*(s)=\arg\max_a Q^*(s,a)
			}
			\]
			
			\vspace{0.4cm}
			
			Другими словами, мы выбираем действие, которое обещает наибольший
			ожидаемый будущий возврат.
			
			
		\end{flushleft}
	\end{frame}
	
	
	
	\begin{frame}{Временное разностное обновление}
		\begin{flushleft}
			
			Вспомним уравнение Беллмана для оценки политики:
			\[
			Q^\pi(s,a)=
			\mathbb{E}_\pi\left[
			r+\gamma Q^\pi(s',a')
			\right].
			\]
			
			\vspace{0.2cm}
			
			Мы не знаем математического ожидания. Вместо этого мы можем использовать
			наблюдаемый переход $(s,a,r,s',a')$.
			
			\vspace{0.2cm}
			
			Мы обновляем нашу оценку функции ценности действия:
			\[
			\boxed{
				Q(s,a)\leftarrow Q(s,a)+\alpha
				\underbrace{
					\left[
					r+\gamma Q(s',a')-Q(s,a)
					\right]
				}_{\text{TD error}}
			}
			\]
			
			где $\alpha$ — скорость обучения.
			
			\medskip
			
			Это называется SARSA-обновлением (State → Action → Reward → State → Action). Это \emph{on-policy} метод: $a'$ определяется политикой.
			
		\end{flushleft}
	\end{frame}
	
	
	\begin{frame}{Q-обучение}
		\begin{flushleft}
			
			Вспомним уравнение оптимальности:
			\[
			Q^*(s,a)=
			\mathbb{E}\left[
			r+\gamma\max_{a'}Q^*(s',a')
			\right].
			\]
			
			\vspace{0.2cm}
			
			Мы можем использовать наблюдаемый переход для обновления нашей оценки:
			\[
			\boxed{
				Q(s,a)\leftarrow Q(s,a)+\alpha
				\left[
				r+\gamma\max_{a'}Q(s',a')
				-Q(s,a)
				\right]
			}
			\]
			
			\vspace{0.3cm}
			
			Это \textbf{Q-обучение}. Это \emph{off-policy} метод — мы используем максимум по всем действиям, чтобы найти $a'$, а не политику.
			
			\vspace{0.2cm}
			
			\begin{itemize}
				\item Мы наблюдаем переход $(s,a,r,s')$.
				\item Мы оцениваем наилучший будущий возврат, используя
				$\max_{a'}Q(s',a')$.
				\item Мы обновляем нашу текущую оценку $Q(s,a)$.
			\end{itemize}
			
		\end{flushleft}
	\end{frame}
	
	
	\begin{frame}{Глубокое Q-обучение}
		\begin{flushleft}
			
			Во многих задачах представление $Q(s,a)$ в виде таблицы
			неосуществимо или нежелательно.
			
			\vspace{0.3cm}
			
			Вместо этого мы можем аппроксимировать функцию ценности действия
			с помощью искусственной нейронной сети (ANN):
			\[
			Q(s,a) \approx Q_\theta(s,a),
			\]
			где $\theta$ обозначает параметры сети.
			
			\vspace{0.3cm}
			
			\begin{itemize}
				\item ANN действует как \textbf{критик}: она оценивает
				ожидаемый будущий возврат для каждого действия.
				\item Вместо обновления отдельных $Q$-значений,
				мы обновляем параметры $\theta$ сети.
				\item Этот подход называется
				\textbf{глубоким Q-обучением} или
				\textbf{Deep Q-Network (DQN)}.
			\end{itemize}
			
		\end{flushleft}
	\end{frame}
	
	
	\begin{frame}{Литература}
		\begin{itemize}
			
			\item Mnih et al. (2013):
			\bref{https://arxiv.org/abs/1312.5602}
			{Playing Atari with Deep Reinforcement Learning. Original Deep Q-Network (DQN) paper.}
			
			\item Hugging Face:
			\bref{https://huggingface.co/learn/deep-rl-course/en/unit2/introduction}
			{Deep Reinforcement Learning Course, Unit 2: Introduction to Q-Learning.}
			
			\item Hugging Face:
			\bref{https://huggingface.co/learn/deep-rl-course/en/unit3/introduction}
			{Deep Reinforcement Learning Course, Unit 3: Deep Q-Learning and DQN.}
			
			\item PyTorch documentation:
			\bref{https://docs.pytorch.org/tutorials/intermediate/reinforcement_q_learning.html}
			{Reinforcement Learning (DQN) Tutorial: implementing DQN on the CartPole environment.}
			
			\item Stable-Baselines3:
			\bref{https://stable-baselines3.readthedocs.io/en/master/modules/dqn.html}
			{DQN documentation and reference implementation.}
			
		\end{itemize}
	\end{frame}
	
	
	\myqrframe
	
	
	
	
	
	
\end{document}

